\documentclass[sn-mathphys-num]{sn-jnl}

\begin{document}
\title[Oxytrees]{Oxytrees: model trees for bipartite learning}
1 -  Keywords

Interaction prediction, model trees, positive-unlabeled learning, Kronecker regularized least squares,  

2 -  Journal
Machine Learning


3 -  3. What is the main claim of the paper? Why is this an important contribution to the machine learning literature?  
In our study, we claim that ensembles of our proposed biclustering tree, the Oxytree, are a state-of-the-art method for bipartite learning, while being substantially faster than previous biclustering forests (which are powerful methods in the field).
We define bipartite learning as a general term for interaction prediction problems in attributed bipartite networks.

Few methods in the literature are designed to predict interactions between unseen instances in a bipartite problem. Furthermore, those who do are either application specific (as often seen in deep learning methods), require costly adaptations (such as global single output methods), or must be partially re-trained (as seen in local methods). A promising alternative are biclustering forests, since they are highly-competitive, applicable to any interaction problem, and do not require hyperparameter tuning. However, traditional biclustering forests pose significant scalability limitations and struggle to model linear relationships in sparse regions of the input space.
Our work addresses these limitations by proposing Oxytrees, highly efficient biclustering model trees that outperform previous biclustering trees while training and predicting up to log(n) times faster, being n the number of training pairs.
%
In conclusion, Oxytrees bring powerful and versatile biclustering forests to a new scale of bipartite learning problems.


4 - What is the evidence you provide to support your claim? Be precise

To the best of our knowledge, we present the largest comparison to date on inductive bipartite interaction prediction. Precisely, we performed cross-validation on 15 datasets and 7 baseline methods, selecting prominent methods from the literature that are able to predict interactions for unseen instances. Furthermore, we also performed experiments in which a percentage of the positive annotations was masked (25\%, 50\% and 75\%), to better investigate the positive-unlabeled setting. Our results highlight ensembles of Oxytrees as a state-of-the-art method, outperforming the baselines in almost all the inductive settings.
We also performed an ablation study and confirm the effectiveness of adding the leaf model to biclustering trees.
Moreover, the computational complexity of Oxytrees in comparison to previous biclustering trees is verified both theoretically and empirically, inducing trees on random data of up to 10^7 interactions.
%The study did not show significant differences in adding a pre-training step with interaction matrix reconstruction, as proposed in the earlier literature.

5 - What papers by other authors make the most closely related contributions, and how is your paper related to them? * 

The studies of Pliakos et al. [1-4] constitute the main literature on biclustering forests. As further discussed in the manuscript, we improved on them, both in terms of performance and computational complexity, by employing model trees, proposing novel induction and inference algorithms, and using a more extensive evaluation setup.

Neighborhood Regularized Logistic Matrix factorization [5] adapts matrix factorization to inductive bipartite learning, and was evaluated as possible component of our method. NRLMF was previously applied to several types of interaction prediction problems.

Regularized least squares with the Kronecker product kernel (RLS-Kron) [6] is a component of our method and an overall strong comparison method in interaction prediction.

[1] Pliakos, Konstantinos, and Celine Vens. "Drug-target interaction prediction with tree-ensemble learning and output space reconstruction." BMC bioinformatics 21 (2020): 1-11.
[2] Pliakos, Konstantinos, Celine Vens, and Grigorios Tsoumakas. ``Predicting drug-target interactions with multi-label classification and label partitioning." IEEE/ACM transactions on computational biology and bioinformatics 18.4 (2019): 1596-1607.
[3] Pliakos, Konstantinos, and Celine Vens. ``Network inference with ensembles of bi-clustering trees." BMC bioinformatics 20 (2019): 1-12.
[4] Pliakos, Konstantinos, Pierre Geurts, and Celine Vens. ``Global multi-output decision trees for interaction prediction." Machine Learning 107 (2018): 1257-1281.
[5] Liu, Yong, et al. ``Neighborhood regularized logistic matrix factorization for drug-target interaction prediction." PLoS computational biology 12.2 (2016): e1004760.
[6] Van Laarhoven, Twan, Sander B. Nabuurs, and Elena Marchiori. ``Gaussian interaction profile kernels for predicting drug–target interaction." Bioinformatics 27.21 (2011): 3036-3043.

6 - Have you published parts of your paper before, for instance in a conference? If so, give details of your previous paper(s) and a precise statement detailing how your paper provides a significant contribution beyond the previous paper(s)
A preliminary version of our work has been published in [6] where the induction algorithm used by our method was proposed. In this manuscript, we extend our previous work in 5 steps: i) we propose model trees which employ regularized least squares with the Kronecker product kernel as their leaf models; ii) a novel significantly faster inference algorithm, which can be used by any decision tree in bipartite learning; iii) a theoretical analysis of the induction algorithm previously published [6]; iv) a more extensive evaluation setup where 5 extra datasets were added, and fractions of the positive annotations were masked;
% v) A novel Python library with Scikit-Learn API for interaction prediction which includes all methods, datasets and evaluation measures employed in this work, allowing reproducible research in the field.

[6] Ilídio, Pedro, André Alves, and Ricardo Cerri. ``Fast Bipartite Forests for Semi-supervised Interaction Prediction." Proceedings of the 39th ACM/SIGAPP Symposium on Applied Computing. 2024.
\end{document}